\chapter{Conclusion}
\label{sec:conclusion}

This thesis set out to make sequential Bayesian experimental design feasible and
effective on stiff, broad-prior dynamical systems where amortized policy methods
fail. It answered the main question in the affirmative, through four results linked
by a single quantity, the likelihood ratio and its prior average $e^{-\Lambda}$.

In answer to RQ1, the prior-contrastive sPCE objective that policy methods are
trained on saturates at $\log(L+1)$ nats, its gradient vanishes at the same rate,
and the contrastive budget needed to escape grows as $e^{\Lambda}$. The failure can
be quantified before training through the escape threshold, which ranges from below
ten to above a million across the suite. In answer to RQ2, drawing contrastives
from a posterior estimator trained once by simulation-based inference removes the
saturation, because the likelihood factor in the importance weight cancels the
divergence that defeated prior contrastives; the cost does not vanish but moves, as
the same posterior concentration reappears in the decline of the effective sample
size. In answer to RQ3, design selection stays valid under that decline, because
the diverging weight scale enters only design-independent terms and cancels in the
ranking, and the criterion approaches a Fisher-information rule in the limit. In
answer to RQ4, PACE matches or beats the policy on every system where the policy
training signal has saturated, including a stiff wastewater model on which the
policy cannot be trained, while learned policies remain stronger on low-information,
high-dimensional designs; the crossover is predicted by the saturation threshold,
not the effective sample size.

The contribution is therefore both a diagnosis and a method. The diagnosis explains
why a widely used training signal fails on an important class of systems, and gives
a cheap test for when it will. The method escapes that failure with a forward-only
estimator that needs neither a trained policy nor gradients through the simulator,
and it remains valid in the regime that makes the problem hard.

The conclusions are qualified by the limitations discussed above. The forward-only
search does not scale to high-dimensional combinatorial designs, PACE returns
reliable rankings rather than calibrated information values under collapsed weights,
and the method assumes a posterior estimator that does not collapse onto the wrong
mode. Future work that addresses these would combine the posterior-contrastive score
with a learned proposal policy to recover scalability, use observation-conditional
inner proposals to keep the effective sample size higher, extend the design rule
beyond a greedy horizon, and test the method on the timed-input biochemical systems
that motivated it.
